% =====================================================================
\chapter{Cài đặt môi trường, hệ điều hành và phần mềm mô phỏng}\label{chap:caidat}
% =====================================================================

\section{Hệ điều hành đã chọn}
\subsection{Trên thiết bị: FreeRTOS cho máy chủ, vòng lặp Arduino cho client}
\begin{table}[H]
\centering\small
\caption{Hệ điều hành trên từng bo mạch}\label{tab:os}
\begin{tabularx}{\linewidth}{@{}lXX@{}}
\toprule
\textbf{Bo mạch} & \textbf{Hệ điều hành / mô hình chạy} & \textbf{Tính năng được dùng} \\ \midrule
ESP32-C3 (máy chủ) & FreeRTOS (bản Espressif, đi kèm ESP-IDF 6.1), tick 100\,Hz & Tác vụ có ưu tiên, mutex, critical section, thông báo tác vụ, hẹn giờ \code{esp\_timer} \\
ESP8266 (client) & Không dùng RTOS: một vòng lặp \code{loop()} không chặn trên Arduino core & Đếm thời gian bằng \code{millis()}, mọi việc chia nhỏ để Wi-Fi luôn được phục vụ \\
\bottomrule
\end{tabularx}
\end{table}

\begin{lydo}[Tại sao máy chủ cần RTOS còn client thì không]
Máy chủ có nhiều việc chạy song song với chu kỳ khác nhau: nhận UDP bất kỳ lúc nào,
phục vụ HTTP, vẽ màn hình, chạy ảnh động 10 khung/giây, đóng thanh chắn đúng giờ.
Viết tất cả trong một vòng lặp buộc phải chia nhỏ mọi việc bằng tay, và chỉ một thao
tác chậm (vẽ màn hình $\approx$ 23\,ms) cũng làm trễ việc khác. FreeRTOS cho phép mỗi
việc là một tác vụ có ưu tiên riêng: việc quan trọng (thanh chắn, nhận bản tin)
luôn chạy trước việc trang trí (ảnh động). Client chỉ làm tuần tự vài việc ngắn
(đọc cảm biến, gửi gói tin), một vòng lặp không chặn là đủ và đơn giản hơn.
\end{lydo}

\subsection{Trên máy phát triển}
Toàn bộ công cụ chạy trên \textbf{Windows 10/11, macOS hoặc Linux}. Nhóm đã phát
triển trên macOS và Windows. Các lệnh trong báo cáo dùng được trên cả ba hệ điều
hành, khác biệt duy nhất là tên cổng nối tiếp (Bảng~\ref{tab:ports}).

\section{Bản quyền phần mềm}
Mọi phần mềm dùng trong đề tài đều \textbf{miễn phí và mã nguồn mở}, nên không cần
bản dùng thử (trial) hay giấy phép trả phí:
\begin{table}[H]
\centering\small
\caption{Phần mềm cần cài và giấy phép}\label{tab:licenses}
\begin{tabularx}{\linewidth}{@{}lXl@{}}
\toprule
\textbf{Phần mềm} & \textbf{Dùng để} & \textbf{Giấy phép} \\ \midrule
ESP-IDF 6.1 (qua ESP-IDF Installation Manager) & Biên dịch, nạp firmware máy chủ & Apache 2.0 \\
Visual Studio Code + ESP-IDF extension (tùy chọn) & Soạn thảo, build/flash bằng giao diện & MIT / Apache 2.0 \\
Arduino IDE 2.x & Biên dịch, nạp hai client & AGPL-3.0 \\
Arduino core ESP8266 & Hỗ trợ bo ESP8266 & LGPL-2.1 \\
MFRC522, DHT sensor library & Thư viện thiết bị & Unlicense / MIT \\
Python 3 & Chạy công cụ giả lập client & PSF \\
Trình biên dịch C/C++ (clang hoặc gcc) & Chạy bộ kiểm thử trên máy tính & Apache 2.0 / GPL \\
Git & Lấy mã nguồn & GPL-2.0 \\
\bottomrule
\end{tabularx}
\end{table}

\section{Cài đặt từng phần mềm}
\subsection{ESP-IDF 6.1 (cho máy chủ)}
\begin{enumerate}
  \item Tải \textbf{ESP-IDF Installation Manager (EIM)} tại
        \url{https://dl.espressif.com/dl/eim/} và cài như ứng dụng thông thường.
  \item Mở EIM, chọn \textbf{ESP-IDF v6.1}, mục tiêu \textbf{esp32c3}, bấm cài đặt.
        EIM tự tải trình biên dịch RISC-V, Python và các công cụ.
  \item Mở terminal ESP-IDF. Windows: shortcut ``ESP-IDF 6.1 PowerShell'' do EIM tạo.
        macOS/Linux: nạp script kích hoạt do EIM tạo, ví dụ:
\begin{lstlisting}[style=shell]
source ~/.espressif/tools/activate_idf_v6.1.sh
\end{lstlisting}
  \item Kiểm tra: \code{idf.py --version} phải in \code{ESP-IDF v6.1}.
\end{enumerate}
ESP32-C3 SuperMini dùng USB Serial/JTAG tích hợp nên \textbf{không cần cài driver}.

\subsection{Arduino IDE và ESP8266 (cho hai client)}
\begin{enumerate}
  \item Cài Arduino IDE 2.x từ \url{https://www.arduino.cc/en/software}.
  \item \emph{File → Preferences → Additional boards manager URLs}, thêm:\\
        \url{https://arduino.esp8266.com/stable/package_esp8266com_index.json}
  \item \emph{Tools → Board → Boards Manager}, tìm \textbf{esp8266}, cài.
  \item \emph{Tools → Manage Libraries}, cài \textbf{MFRC522} (GithubCommunity) và
        \textbf{DHT sensor library} (Adafruit, chấp nhận cài kèm Adafruit Unified Sensor).
  \item Bo NodeMCU dùng chip USB-UART CH340 hoặc CP2102: nếu máy không nhận cổng,
        cài driver tương ứng của nhà sản xuất chip.
\end{enumerate}

\subsection{Công cụ cho giả lập và kiểm thử}
\begin{itemize}
  \item \textbf{Python 3.8+}: có sẵn trên macOS/Linux; Windows tải tại python.org. Không cần thư viện ngoài.
  \item \textbf{Trình biên dịch C/C++}: macOS \code{xcode-select --install}; Ubuntu
        \code{sudo apt install build-essential}; Windows dùng MSYS2/MinGW hoặc WSL.
\end{itemize}

\section{Phần mềm mô phỏng và các tính năng chính}
Đề tài dùng hai công cụ mô phỏng do nhóm tự viết, giúp phát triển và kiểm tra hệ
thống khi chưa có (hoặc không mang theo) đủ phần cứng.

\subsection{Bộ giả lập client bằng Python (\file{examples/})}
\begin{table}[H]
\centering\small
\caption{Công cụ giả lập client}\label{tab:simulators}
\begin{tabularx}{\linewidth}{@{}lX@{}}
\toprule
\textbf{Công cụ} & \textbf{Tính năng} \\ \midrule
\file{udp\_sender.py} & Giả lập Sensor Client: gửi một bản tin \code{parking\_update} (5 ô đỗ, mực nước, cháy) tới máy chủ thật qua UDP broadcast \\
\file{udp\_rfid\_sender.py} & Giả lập RFID Client: gửi \code{rfid\_event} với UID bất kỳ và cổng \code{--gate in|out}; chạy hai lần để giả lập vào rồi ra \\
\file{udp\_sender\_espidf.c} & Ví dụ gửi bản tin từ một bo ESP-IDF khác \\
\bottomrule
\end{tabularx}
\end{table}
Máy tính chỉ cần kết nối Wi-Fi \code{ESP32-Parking}; máy chủ không phân biệt bản tin
đến từ ESP8266 thật hay từ script, nên toàn bộ màn hình, thanh chắn, còi, trang web và
nhật ký đều được kiểm tra như thật.

\subsection{Mô phỏng phần cứng Arduino trên máy tính (\file{tests/arduino/})}
Đây là lớp giả lập các thư viện Arduino/ESP8266: đồng hồ \code{millis()} điều khiển
được, mức logic của từng chân, giá trị ADC, Wi-Fi bật/tắt, gói UDP được ghi lại,
đầu đọc RFID có thể ``đặt thẻ'' và DHT11 trả nhiệt độ tùy ý. Nhờ vậy \emph{chính mã
nguồn} \file{SensorClient.ino} và \file{RFIDClient.ino} được biên dịch và chạy trên
máy tính, mỗi bản tin chúng gửi được đưa thẳng vào bộ đọc JSON thật của máy chủ để
kiểm tra hai bên ``nói cùng một ngôn ngữ''. Chi tiết ở Chương~\ref{chap:kiemthu}.

\begin{lydo}[Tại sao tự viết công cụ mô phỏng thay vì dùng trình mô phỏng có sẵn]
Các trình mô phỏng vi điều khiển trực tuyến hỗ trợ ESP32-C3 nhưng không mô phỏng
được việc ba bo mạch nói chuyện với nhau qua Wi-Fi, và không có mô hình cho cả
MFRC522 lẫn module lửa/nước cụ thể của nhóm. Giả lập ở mức \emph{giao thức} (script
Python gửi bản tin thật) và ở mức \emph{thư viện} (đồng hồ, chân, UDP giả) kiểm tra
đúng phần logic nhóm tự viết, chạy trong vài giây, lặp lại được và không phụ thuộc
dịch vụ bên ngoài.
\end{lydo}
